The analysis followed four connected stages. First, the MLCW and cGNSS displacement records were evaluated on common monthly dates. Second, the sediment proportions were expressed as independent logratio coordinates. Third, the monthly observations from all six depth sections were organized for station-specific regression. Finally, the fitted model was tested in temporal order while consecutive MLCW observations remained unavailable for six months. This sequence is described below.

\subsection{Deformation time series model}
\label{subsec:deformation_model}

The MLCW and cGNSS records did not share a uniform set of observation dates. Direct differencing of their raw measurements would therefore compare intervals with unequal starting and ending dates. A deformation time series model was fitted to each displacement series so that all series could be evaluated on the same monthly dates.

The model represented displacement as the sum of a reference value, a linear component, annual and semiannual components, and offsets at documented discontinuities. For observation time $t$, the fitted displacement was
\begin{equation}
d(t)=d_0+v(t-t_0)
+\sum_{k\in\{1,2\}}
\left[a_k\sin(2\pi k t)+b_k\cos(2\pi k t)\right]
+\sum_{j=1}^{J}c_j H(t-t_j),
\label{eq:deformation_model}
\end{equation}
where $d_0$ denotes displacement at reference time $t_0$, $v$ denotes linear velocity, $a_k$ and $b_k$ describe the periodic components, and $H(t-t_j)$ represents a step at discontinuity time $t_j$ \citep{nikolaidis_observation_2002}. Evaluation of each fitted series at the common monthly dates produced aligned cumulative displacement records. Differences between consecutive fitted values then represented the MLCW compaction and cGNSS vertical displacement that occurred during each month.

% ===== PROPOSED CORRECTION -- AWAITING AUTHOR APPROVAL -- DO NOT UNCOMMENT WITHOUT REVIEW =====
% ISSUE: the last sentence of this paragraph ("Differences between consecutive fitted values
% then represented the MLCW compaction and cGNSS vertical displacement that occurred during
% each month.") is correct for cGNSS but incorrect for MLCW. Verified against
% 23_build_input_snapshot.py (load_mlcw reads the raw CSV directly; y_observed = mlcw.diff(),
% "ZERO outlier-related modification" per that script's own docstring) and independently
% confirmed in run048_tuku_no_update_sensitivity.py's handoff docstring (lines 893-898):
% fitted-value differencing is true for cGNSS, false for MLCW.
%
% PROPOSED REPLACEMENT for the final sentence of this paragraph:
% Differences between consecutive fitted cGNSS values then represented the vertical
% displacement that occurred during each month. The MLCW compaction increment for the same
% month was instead calculated as the first difference of the corresponding raw monthly
% MLCW record, without passing through the fitted deformation model, because the MLCW record
% did not require alignment to a different observation calendar.
%
% AUTHOR DECISION NEEDED (not resolved by this proposal): the opening sentence of this
% subsection, line 6 ("A deformation time series model was fitted to each displacement
% series so that all series could be evaluated on the same monthly dates"), says the model
% was fitted to EACH/ALL series. If MLCW's monthly increment in the analysis pipeline is in
% fact raw-differenced (as verified above), "each/all" in line 6 is broader than what the
% pipeline does for MLCW, and the paragraph will read as internally inconsistent once the
% line-16 correction above is adopted. Please confirm whether line 6 describes a separate
% geodetic/display alignment product (in which case no change is needed there) or whether it
% must also be narrowed to "cGNSS" only.
% ===== END PROPOSED CORRECTION =====

\subsection{Isometric logratio transformation of sediment composition}
\label{subsec:ilr_transformation}

Sediment composition provided time-independent information that distinguished the six depth sections. Each section contained four proportions, gravel, coarse sand, fine sand, and fine-grained deposits, whose sum was fixed at 100\%. An increase in one proportion therefore required a decrease in at least one other proportion. Direct use of all four values would incorrectly treat them as independent predictors.

An isometric logratio transformation represented the relative composition with three independent coordinates \citep{egozcue_isometric_2003, oh_using_2024}. A sequential binary partition first divided the four sediment components into physically interpretable groups. Each coordinate then expressed the logarithmic ratio between the geometric means of two groups. For groups containing $r$ and $s$ components, the corresponding balance was
\begin{equation}
z=\sqrt{\frac{rs}{r+s}}
\ln\left(\frac{g(\boldsymbol{x}_{+})}{g(\boldsymbol{x}_{-})}\right),
\label{eq:ilr_balance}
\end{equation}
where $g(\boldsymbol{x}_{+})$ and $g(\boldsymbol{x}_{-})$ denote the geometric means of the components in the two groups. The three balances retained the lithological differences among the Tuku sections without assigning unmeasured hydraulic or mechanical properties to the logged sediment types.

\subsection{Preparation of model inputs}
\label{subsec:predictor_preparation}

Each model observation represented one depth section during one month. Its response was the MLCW compaction increment for that section. The corresponding predictors described three types of information available during the same month. Hydraulic head changes represented conditions within the aquifer system, cGNSS displacement represented the integrated surface response, and the sediment balances identified the lithological setting of the section.

The hydraulic predictors included information associated with the section being estimated and concurrent observations from the other Tuku sections. Earlier hydraulic head and cGNSS changes were also retained where they represented a delayed compaction response. Annual and semiannual terms described recurring seasonal variation. Section indicators and selected interactions allowed the relation between hydraulic head change and compaction to differ with depth, season, and sediment composition. \Cref{tab:predictor_summary} summarizes the physical role of each predictor group. The final variable names, temporal transformations, and lag periods will be listed in \Cref{app:predictors} after the Tuku analysis is frozen.

Observations from all six sections formed one station-specific calibration dataset. This pooled structure allowed sections with longer or clearer records to contribute to the estimation of common coefficients while preserving section-dependent relations. Predictor centering and scaling were calculated from the calibration period available before each evaluation block. The same values were then applied to the following six evaluation months. MLCW compaction from earlier months was not included as a predictor, so every estimate depended on the auxiliary monitoring records rather than on a previous compaction value.

\subsection{Bayesian ridge regression}
\label{subsec:bayesian_ridge}

Bayesian ridge regression was used because the single-station dataset contained several related hydraulic, seasonal, and section-dependent predictors. Ordinary least squares can assign unstable coefficients when predictors contain overlapping information. Bayesian ridge regression reduces this instability by shrinking weakly supported coefficients toward zero while retaining a linear relation that can be inspected by predictor group.

For $n$ monthly section observations and $p$ predictors, the model was
\begin{equation}
\boldsymbol{y}=\boldsymbol{X}\boldsymbol{\beta}+\boldsymbol{\varepsilon},
\qquad
\boldsymbol{\varepsilon}\sim\mathcal{N}(\boldsymbol{0},\alpha^{-1}\boldsymbol{I}),
\label{eq:brr_likelihood}
\end{equation}
where $\boldsymbol{y}$ contains monthly section compaction, $\boldsymbol{X}$ contains the standardized predictors, $\boldsymbol{\beta}$ contains their coefficients, and $\alpha$ denotes the error precision. A zero-centered Gaussian prior regularized the coefficients,
\begin{equation}
\boldsymbol{\beta}\sim\mathcal{N}(\boldsymbol{0},\lambda^{-1}\boldsymbol{I}),
\label{eq:brr_prior}
\end{equation}
where $\lambda$ controls the amount of shrinkage. The posterior mean of the fitted response provided the monthly compaction estimate. The regression model described statistical associations in the monitoring record and was not interpreted as a substitute for a coupled groundwater flow and aquifer-system compaction model.

\subsection{Evaluation under a six-month data delay}
\label{subsec:delayed_evaluation}

The evaluation reproduced a monitoring schedule in which MLCW measurements continued each month but were delivered in batches after six months. An initial period from \placeholder{INITIAL CALIBRATION START} to \placeholder{INITIAL CALIBRATION END} provided the first calibration dataset. The fitted model then estimated compaction for each of the next six months from GWL and cGNSS observations available during that month. The six corresponding MLCW responses remained hidden throughout the block.

All six hidden MLCW observations became available at the end of the block. Their release allowed the estimates to be evaluated and the observed responses to be added to the calibration dataset. The model was then fitted again before estimation began for the next block. Nonoverlapping six-month blocks continued through \placeholder{FINAL EVALUATION MONTH}. This expanding calibration window preserved temporal order because no response from the current or a future block could influence its estimates. \Cref{fig:delayed_observation_design} illustrates one update cycle.

Three reference estimates placed the updated model in context. A frozen model retained the coefficients fitted during the initial calibration period and therefore showed the value of periodic updating. A last-observation baseline repeated the most recent monthly compaction value available before each block. A seasonal baseline repeated the compaction observed in the corresponding month of the preceding year. Both persistence baselines used only MLCW observations that would have been available when the block began.

Performance was calculated for each depth section and for all sections combined. The coefficient of determination ($R^2$) described agreement in temporal variation, while root mean square error (RMSE) and mean absolute error (MAE) measured the magnitude of estimation errors in \placeholder{CONFIRM ERROR UNIT}. Performance was also grouped by the first through sixth month after each model update. This comparison tested whether the estimates became less reliable as the time since the latest MLCW batch increased.

\subsection{Prediction intervals and uncertainty evaluation}
\label{subsec:uncertainty}

A point estimate alone does not show the range of monthly compaction values that remains consistent with past model errors. Error-calibrated prediction intervals were therefore constructed around each estimate. The interval width was based only on errors from evaluation blocks that had already been completed, which prevented observations from the current or future block from informing its uncertainty.

Absolute errors were grouped by depth section and by the number of months since the latest model update. Before block $b$, the eligible errors for section $s$ and update distance $h$ defined a finite-sample 90th percentile, denoted $q_{b,s,h}$ \citep{angelopoulos_conformal_2023}. The prediction interval around estimate $\widehat{y}_{b,s,h}$ was
\begin{equation}
\left[
\widehat{y}_{b,s,h}-q_{b,s,h},
\widehat{y}_{b,s,h}+q_{b,s,h}
\right].
\label{eq:prediction_interval}
\end{equation}
An interval was reported only after at least \placeholder{MINIMUM ELIGIBLE ERROR COUNT} earlier errors were available for the corresponding group. \Cref{app:intervals} will document the final eligibility and quantile rules after verification against the operational code.

Empirical coverage was the proportion of observed monthly compaction values enclosed by the reported intervals. Mean interval width described how broad those intervals were. Coverage and width were examined together because a wider interval can increase coverage without providing a precise estimate \citep{singh_uncertainty_2024}. The nominal 90\% level was treated as a calibration target rather than a guarantee because the monitoring record contains temporal dependence and can change through time \citep{zaffran_adaptive_2022}.

% ===== PROPOSED ADDITION -- AWAITING AUTHOR APPROVAL -- DO NOT UNCOMMENT WITHOUT REVIEW =====
% NOTE TO AUTHOR: this subsection documents a separate evaluation design (no refit within a
% horizon) from subsec:delayed_evaluation above (which refits every six months and is left
% unchanged by this proposal). The two are not interchangeable and should not share a label,
% a distance-since-update grouping, or the q_{b,s,h} conformal formula from subsec:uncertainty.
% The source handoff document for this analysis (run048_tuku_no_update_sensitivity.py, module
% docstring and "Evaluation design" / "Terminology" sections) suggests this procedure could
% instead be sited as a Discussion methods paragraph or supplement rather than in Methods
% proper -- please confirm placement before activating.
%
% \subsection{No-update reduced-observation sensitivity analysis}
% \label{subsec:no_update_sensitivity}
%
% A separate analysis evaluated model reliability under a hypothetical schedule in which the
% MLCW compaction well was not checked between scheduled readings, in contrast to the periodic
% refitting described in \Cref{subsec:delayed_evaluation}. Rolling origins were defined at the
% start of each existing six-month evaluation block used in \Cref{subsec:delayed_evaluation},
% giving 22 origins spaced six months apart from May 2013 through November 2023. At each
% origin, the Bayesian ridge regression model of \Cref{subsec:bayesian_ridge} was fitted once,
% using only monthly GWL and cGNSS observations and the MLCW history available before that
% origin, and was not refitted for the following twelve months. Estimates for the twelve
% months after each origin were produced from this single fit using the monthly GWL and cGNSS
% observations that would have been available during the corresponding month, without access
% to MLCW observations from within that twelve-month horizon. The first six months of each
% twelve-month estimate series were also reported separately, as a six-month horizon derived
% from the same fitted model and the same origin rather than from an independently refitted
% model. Consecutive twelve-month horizons therefore overlapped by six months. A twenty-third
% origin, 1 May 2024, had only six subsequent months of monitoring data available and could
% not support a twelve-month horizon; it was evaluated separately as a six-month diagnostic
% case and excluded from the 22-origin comparison.
%
% Estimation error was calculated for each depth section and for all sections combined, at
% both the six-month and twelve-month horizons, as the cumulative difference between the
% summed monthly estimates and the summed monthly MLCW observations released at the end of
% each horizon, once those observations became available. Cumulative error was signed as the
% cumulative estimate minus the cumulative observation, so that a negative value indicated
% underestimated compaction magnitude when the observed cumulative change was itself negative.
% Cumulative absolute error and cumulative signed error (bias) were reported in millimeters.
% A horizon-normalized error rate was also calculated by dividing the cumulative absolute
% error by the number of months in the horizon, reported in millimeters per month. The
% coefficient of determination was not calculated for this analysis because no continuous
% monthly reference series was withheld for comparison within a horizon.
%
% Ninety-five percent confidence intervals for each error summary, and for the difference
% between the twelve-month and six-month horizon-normalized error rates, were constructed with
% a paired moving-block bootstrap over the ordered sequence of 22 origins, using a block length
% of two origins and 2000 resamples, applied identically and in lockstep to the six-month and
% twelve-month results so that each resample compared the same subset of origins at both
% horizons.
% ===== END PROPOSED ADDITION =====

\begin{figure}[htbp]
  \centering
  \resizebox{\linewidth}{!}{%
  \begin{tikzpicture}[
    node distance=0.65cm,
    box/.style={draw, rectangle, rounded corners=2pt, align=center, minimum width=3.35cm, minimum height=1.05cm, font=\small},
    arrow/.style={-{Stealth[scale=0.9]}, line width=0.8pt}
  ]
    \node (train) [box, fill=blue!7] {Available history\\Fit the model};
    \node (estimate) [box, fill=orange!8, right=of train] {Months 1--6\\Estimate with current GWL and cGNSS};
    \node (release) [box, fill=green!8, right=of estimate] {End of month 6\\Receive six monthly MLCW records};
    \node (update) [box, fill=gray!10, below=of release] {Compare, append records,\\and update the model};
    \draw [arrow] (train) -- (estimate);
    \draw [arrow] (estimate) -- (release);
    \draw [arrow] (release) -- (update);
    \draw [arrow] (update.west) -| (train.south);
  \end{tikzpicture}%
  }
  \caption{Evaluation of monthly compaction estimates under a six-month delay in MLCW data delivery. Each block was estimated from monitoring information available during its six target months. The corresponding MLCW observations were released together at the end of the block and were then used to evaluate and update the model.}
  \label{fig:delayed_observation_design}
\end{figure}

\begin{table}[htbp]
  \centering
  \small
  \renewcommand{\arraystretch}{1.15}
  \caption{Predictor groups used to estimate monthly compaction at Tuku.}
  \label{tab:predictor_summary}
  \begin{tabularx}{\textwidth}{@{}p{0.21\textwidth}p{0.27\textwidth}X@{}}
    \toprule
    \textbf{Predictor group} & \textbf{Information represented} & \textbf{Role in the model} \\
    \midrule
    cGNSS displacement & Current and earlier surface displacement & Described the integrated vertical response at Tuku \\
    Target-section hydraulic head & Current and earlier hydraulic head changes & Described hydraulic conditions associated with the section being estimated \\
    Other-section hydraulic head & Concurrent head changes in the remaining sections & Described hydraulic conditions elsewhere in the monitored profile \\
    Seasonal terms & Annual and semiannual variation & Represented recurring seasonal patterns \\
    Sediment composition & Three isometric logratio balances & Distinguished material composition among depth sections \\
    Section and interaction terms & Relations that vary with depth, season, and composition & Allowed one station-specific model to represent different section responses \\
    \bottomrule
  \end{tabularx}
\end{table}

\FloatBarrier
